\begin{textsample}{POS Dim 1 – vem\_ed\_09 – Score 24.00 – vem\_ed\_09/t035.txt}  \label{ex:f1_pos_083}
QUEM \textbf{É} QUEM Victoria Traverso Castro \textbf{é} mestre em Linguística e assessora de internacionalização de alunos, ex-alunos e docentes da Inacap (Chile).
Segue seu depoimento sobre a internacionalização da Inacap.
O trabalho em conjunto com as áreas acadêmicas e unidades da Inacap tem \textbf{sido} fundamental para o êxito da internacionalização em \textbf{casa}.
Temos projetos \textbf{curriculares} e extracurriculares e ambos incluem um trabalho árduo com professores, diretores e assessores em \textbf{cada} unidade e nas áreas acadêmicas centralizadas.
Sem o \textbf{apoio} desses atores \textbf{relevantes}, não \textbf{seria} \textbf{possível} \textbf{implementar} os \textbf{programas} de Internacionalização em Casa.
Dos projetos \textbf{curriculares}, temos entre 10 e 13 programados para este semestre.
Eles têm um alcance \textbf{amplo}, pois impactam \textbf{todos} os estudantes matriculados na \textbf{disciplina}.
Por exemplo, na área de negócios, esperamos envolver cerca de 5 mil estudantes, e apesar de \textbf{que} essas atividades tão massivas geralmente não \textbf{sejam} COILs [ Collaborative Online International Learning ], \textbf{representam} Internacionalização em Casa.
Para este semestre temos programados 8 COILs propriamente ditos, vários deles com o Centro Paula Souza e muitos \textbf{são} \textbf{continuações} de projetos \textbf{iniciados} no semestre \textbf{anterior} e \textbf{que} tiveram muito sucesso.
Trabalhar com CPS / Cesu \textbf{é} sempre um prazer.
Não somente os brasileiros \textbf{são} muito simpáticos, mas também trabalham duro e estão sempre dispostos a colaborar, o \textbf{que} facilita muitíssimo o trabalho.
Tanto Osvaldo [ Succi Junior ] como os outros professores têm \textbf{sido} \textbf{essenciais} no \textbf{desenvolvimento} de atividades de Internacionalização em Casa.
Aprendemos muitíssimo com \textbf{cada} um deles.
As atividades COIL \textbf{começaram} durante 2020 e o nosso primeiro colaborador \textbf{foi} CPS / Cesu.
Em 2020 houve 2 projetos e no primeiro semestre de 2021, \textbf{foram} 9 projetos e algumas atividades de internacionalização, como \textbf{encontros} estudantis, além das 14 atividades \textbf{curriculares}, \textbf{que} alcançaram mais de 4 mil estudantes.
Neste semestre há um \textbf{número} similar de projetos em andamento.
Como podem ver, temos muito trabalho pela frente!
Nosso objetivo central para os próximos anos \textbf{é} \textbf{que} a instituição compreenda a \textbf{importância} da \textbf{interculturalidade} e internacionalização de nossos estudantes.
Para isso, trabalhamos arduamente para massificar e democratizar as experiências internacionais, fazendo o \textbf{possível} para \textbf{que} atinjam \textbf{cada} vez mais estudantes e docentes.
Esperamos \textbf{que}, por meio dessas experiências, nossos estudantes melhorem sua empregabilidade, \textbf{sejam} melhores cidadãos globais e busquem trabalhar pelo bem \textbf{comum} em sua vida \textbf{profissional}, integrando os Objetivos de Desenvolvimento \textbf{Sustentável} da ONU e comunicando-se de maneira efetiva em vários \textbf{contextos} e com \textbf{pessoas} diversas.

% matched lemmas: amplo, anterior, apoio, cada, casa, começar, comum, contexto, continuação, curricular, desenvolvimento, disciplina, encontro, essencial, implementar, importância, iniciar, interculturalidade, número, pessoa, possível, profissional, programa, que, relevante, representar, ser, sustentável, todo
\end{textsample}
